\mychapter{0}{Résumé}


\par L'architecture microservice propose de construire une application en tant qu'ensemble de services autonomes et indépendants. Chaque service est conçu pour se concentrer sur un module spécifique et peut communiquer avec les autres services en cas de besoin. Cette approche offre plusieurs avantages tels que la scalabilité, la flexibilité et la résilience du système. 

\par Les microservices sont devenus très populaires vu qu'ils facilitent la gestion et la maintenance des systèmes informatiques.Chaque service peut être développé, déployé et mis en échelle de manière indépendante ce qui aide à l'évolution correcte des applications.


\par Dans le cadre de mon projet de fin d’études pour l’obtention du diplôme d'ingénieur d'état à l’École Marocaine des Sciences de l'Ingénieur, réalisé au sein de l’entreprise SQLI Digital Experience, j’ai intégré une équipe ambitieuse qui travaille sur la migration du projet e-commerce de Nespresso vers une architecture microservice tout en se basant  sur les notions du "Clean Code" et du "Domain Driven Design". Ce projet vise à séparer le module "panier" de l'ancienne architecture et le développer avec ces nouvelles approches.


\par Pour réaliser le projet, nous avons tout d’abord analysé l'existant pour définir ses limites ainsi que les exigences du module à développer au début. Ensuite, nous avons proposé une nouvelle architecture et élaboré une conception en se basant sur plusieurs diagrammes, et enfin, nous avons procédé à la réalisation du projet en utilisant les différentes technologies décrites dans le rapport.

\vspace{1cm}


\noindent\rule[2pt]{\textwidth}{0.5pt}

{\textbf{Mots clés :}}
microservice, architecture onion, ECAPI split, migration, e-commerce


\noindent\rule[2pt]{\textwidth}{0.5pt}

\clearpage

\mychapter{0} 
The microservice architecture proposes to build an application as a collection of autonomous and independent services. Each service is designed to focus on a specific module and can communicate with other services when needed. This approach offers several advantages such as scalability, flexibility, and system resilience.

Microservices have become very popular as they facilitate the management and maintenance of IT systems. Each service can be developed, deployed, and scaled independently, which helps ensure the proper evolution of applications.

As part of my final year project to obtain a state engineer degree at the Moroccan School of Engineering Sciences, carried out at SQLI Digital Experience company, I joined an ambitious team working on the migration of the Nespresso e-commerce project to a microservice architecture, based on the principles of "Clean Code" and "Domain-Driven Design". This project aims to separate the "cart" module from the old architecture and develop it using these new approaches.

To accomplish the project, we first analyzed the existing system to identify its limitations and the requirements of the module to be developed initially. Then, we proposed a new architecture and created a design based on many diagrams. Finally, we proceeded with the implementation of the project, using the different technologies described in the report.

\vspace{1cm}



\noindent\rule[2pt]{\textwidth}{0.5pt}

{\textbf{Keywords :}}
microservice, onion architecture, ECAPI split, migration, e-commerce
\\
\noindent\rule[2pt]{\textwidth}{0.5pt}
